\begingroup
\renewcommand{\arraystretch}{1.45}
\InterpolationArticle{22}
Wir gehen zum anderen Fall von Artikel 11 über, in dem die Reihe \(X\) zwar dieselbe bleibt wie in Artikel 20, die Anzahl der gegebenen Werte aber um eins kleiner ist als
die Anzahl der unbekannten Koeffizienten, nämlich \(=2m=\mu\). Hier gilt
\SourceLeaf{308}
\[
\begin{aligned}
&\cos\frac12(t-a)\sin\frac12(t-b)\sin\frac12(t-c)\sin\frac12(t-d)\mathinner{.\,.\,.}\\[5pt]
&=-\frac{\sin\frac12\mu(t-a)\cos\frac12(t-a)}{2^{\mu-1}\sin\frac12(t-a)}\\[5pt]
&=-\frac{\cos\frac12(t-a)}{2^{\mu-2}}\bigl(\cos\frac12(t-a)+\cos\frac32(t-a)\\[5pt]&\qquad{}+\cos\frac52(t-a)+\mathinner{.\,.\,.}+\cos\frac{\mu-1}{2}(t-a)\bigr)\\[5pt]
&=-\frac{1}{2^{\mu-1}}\bigl(1+2\cos(t-a)+2\cos2(t-a)\\[5pt]&\qquad{}+2\cos3(t-a)+\mathinner{.\,.\,.}+2\cos(m-1)(t-a)\\[5pt]
&\qquad +\cos m(t-a)\bigr)
\end{aligned}
\]
In der Formel von Artikel 11 für $T$ wird also der Koeffizient von $A$
\[
\begin{aligned}
&=\frac1\mu\bigl(1+2\cos(t-a)+2\cos2(t-a)\\[5pt]&\qquad{}+2\cos3(t-a)+\mathinner{.\,.\,.}\\[5pt]
&\qquad +2\cos(m-1)(t-a)+\cos m(t-a)\bigr)
\end{aligned}
\]
Ganz ähnliche Ausdrücke (wenn nur $a$ in $b,c,d$ usw. geändert wird) findet man für die Koeffizienten von $B,C,D$ usw. Schließlich wird der letzte Teil von $T$
\[
\begin{aligned}
&2^{2m-1}i^{2m}\sin\frac12(t-a)\sin\frac12(t-b)\sin\frac12(t-c)\sin\frac12(t-d)\mathinner{.\,.\,.}\\[5pt]&\qquad\times\{\alpha^m\cos\frac12s+\beta^m\sin\frac12s\}
\end{aligned}
\]
zu
\[
\begin{aligned}
&=-i^{2m}\sin m(t-a)\{\alpha^m\cos(ma+(2m-1)90^{0})\\[5pt]&\qquad{}+\beta^m\sin(ma+(2m-1)90^{0})\}\\[5pt]
&=-\sin m(t-a)(\alpha^m\sin ma-\beta^m\cos ma)
\end{aligned}
\]
Daraus folgt schließlich
\[
\begin{aligned}
&T=\frac1\mu(A+B+C+D+\mathrm{usw.})\\[5pt]
&\quad+\frac2\mu(A\cos a+B\cos b+C\cos c+D\cos d+\mathrm{usw.})\cos t\\[5pt]
&\quad+\frac2\mu(A\sin a+B\sin b+C\sin c+D\sin d+\mathrm{usw.})\sin t\\[5pt]
&\quad+\frac2\mu(A\cos2a+B\cos2b+C\cos2c+D\cos2d+\mathrm{usw.})\cos2t\\[5pt]
&\quad+\frac2\mu(A\sin2a+B\sin2b+C\sin2c+D\sin2d+\mathrm{usw.})\sin2t\\[5pt]
&\quad+\mathrm{usw.}\\[5pt]
&\quad+\frac1\mu(A\cos ma+B\cos mb+C\cos mc+D\cos md+\mathrm{usw.})\cos mt\\[5pt]
&\quad+\frac1\mu(A\sin ma+B\sin mb+C\sin mc+D\sin md+\mathrm{usw.})\sin mt\\[5pt]
&\quad+(\alpha^m\sin ma-\beta^m\cos ma)\\[5pt]&\qquad(\sin ma\cos mt-\cos ma\sin mt)
\end{aligned}
\]
Stellt man diesen Ausdruck nach Weglassung seines letzten Teils durch
\[
\begin{aligned}
&\gamma+\gamma'\cos t+\gamma''\cos2t+\mathrm{usw.}+\gamma^m\cos mt\\[5pt]
&\quad+\delta'\sin t+\delta''\sin2t+\mathrm{usw.}+\delta^m\sin mt
\end{aligned}
\]
\SourceLeaf{309}
dar, so haben offenbar alle Koeffizienten $\gamma,\gamma',\delta',\gamma'',\delta''$ usw. genau dieselben Werte wie im vorigen Artikel, wo $\mu=2m+1$ war, mit Ausnahme der beiden letzten $\gamma^m,\delta^m$, die dort doppelt so groß waren. Ferner ist klar, dass alle Koeffizienten
$\alpha,\alpha',\alpha''\mathinner{.\,.\,.}\alpha^{m-1},\beta',\beta''\mathinner{.\,.\,.}\beta^{m-1}$ jeweils gleich $\gamma,\gamma',\gamma''\mathinner{.\,.\,.}\gamma^{m-1},\delta',\delta''\mathinner{.\,.\,.}\delta^{m-1}$ und damit bekannt werden, während zwischen den beiden letzten $\alpha^m,\beta^m$ die Gleichungen
\[
\begin{aligned}
&\alpha^m=\gamma^m+(\alpha^m\sin ma-\beta^m\cos ma)\sin ma\\[5pt]
&\beta^m=\delta^m-(\alpha^m\sin ma-\beta^m\cos ma)\cos ma
\end{aligned}
\]
bestehen, die jedoch (wie sich aus der Natur der Aufgabe vorhersehen ließ) zur vollständigen Bestimmung der Koeffizienten $\alpha^m,\beta^m$ nicht ausreichen. Da nämlich
\[
\begin{aligned}
&\cos ma=-\cos mb=\cos mc=-\cos md\ \mathrm{usw.}
\end{aligned}
\]
und
\[
\begin{aligned}
&\sin ma=-\sin mb=\sin mc=-\sin md\ \mathrm{usw.}
\end{aligned}
\]
und daher
\[
\begin{aligned}
&\gamma^m=\frac{\cos ma}{\mu}(A-B+C-D+\mathrm{usw.})\\[5pt]
&\delta^m=\frac{\sin ma}{\mu}(A-B+C-D+\mathrm{usw.})
\end{aligned}
\]
gilt, erkennt man leicht, dass beide Gleichungen in der einzigen Gleichung
\[
\begin{aligned}
&\alpha^m\cos ma+\beta^m\sin ma=\frac1\mu(A-B+C-D+\mathrm{usw.})
\end{aligned}
\]
enthalten sind. Es muss also außerdem eine weitere Beziehung zwischen $\alpha^m$ und $\beta^m$ gegeben sein, damit sich beide Koeffizienten vollständig bestimmen lassen.

\Needspace{36\baselineskip}
\InterpolationArticle{23}
Bricht die Reihe $X$ nicht mit $\cos mx$ und $\sin mx$ ab, so wird der im vorigen Artikel gefundene Wert von $T$
unvollständig sein, und die Koeffizienten $\gamma,\gamma',\delta'$ usw. werden von den wahren Werten
$\alpha,\alpha',\beta'$ usw. abweichen. Man erkennt leicht, dass sich, wenn die Koeffizienten der folgenden Glieder in $X$ mit denselben Zeichen bezeichnet werden, die wir in Artikel 21 verwendet haben,
für $\gamma,\gamma',\delta'$ usw. bis zu $\gamma^{m-1},\delta^{m-1}$ genau dieselben Werte ergeben,
die wir dort hergeleitet haben; die beiden letzten werden dagegen
\[
\begin{aligned}
&\gamma^m=\frac12\alpha^m+\frac12(\alpha^m+\alpha^{3m})\cos\mu a\\[5pt]&\qquad{}+\frac12(\alpha^{3m}+\alpha^{5m})\cos2\mu a+\mathrm{usw.}\\[5pt]
&\qquad +\frac12(\beta^m+\beta^{3m})\sin\mu a+\frac12(\beta^{3m}+\beta^{5m})\sin2\mu a+\mathrm{usw.}\\[5pt]
&\delta^m=\frac12\beta^m-\frac12(\beta^m-\beta^{3m})\cos\mu a\\[5pt]&\qquad{}-\frac12(\beta^{3m}-\beta^{5m})\cos2\mu a-\mathrm{usw.}\\[5pt]
&\qquad -\frac12(\alpha^m-\alpha^{3m})\sin\mu a-\frac12(\alpha^{3m}-\alpha^{5m})\sin2\mu a-\mathrm{usw.}
\end{aligned}
\]
\SourceLeaf{310}
In diesem Fall darf man also die Koeffizienten $\gamma,\gamma',\delta',\gamma'',\delta''\mathinner{.\,.\,.}\gamma^{m-1},\delta^{m-1}$ für die wahren Werte
$\alpha,\alpha',\beta',\alpha'',\beta''\mathinner{.\,.\,.}\alpha^{m-1},\beta^{m-1}$ nehmen, soweit die folgenden Koeffizienten $\alpha^{m+1},\beta^{m+1}$ usw.
vernachlässigt werden können; auf die beiden letzten $\gamma^m,\delta^m$ darf dieser Schluss jedoch nicht ausgedehnt werden, sofern sie nicht selbst als vernachlässigbar angesehen werden können.

\endgroup
